\documentclass[12pt,a4paper]{article}
\usepackage[utf8]{inputenc}
\usepackage[T1]{fontenc}
\usepackage{amsmath}
\usepackage{amsfonts}
\usepackage{amssymb}
\usepackage{graphicx}
\usepackage{geometry}
\usepackage{setspace}
\usepackage{titlesec}
\usepackage{fancyhdr}
\usepackage{hyperref}
\usepackage{booktabs}
\usepackage{array}
\usepackage{multirow}
\usepackage{longtable}
\usepackage{pgfplots}
\usepackage{tikz}
\usepackage[backend=biber,style=authoryear,sorting=nyt]{biblatex}
\usepackage{algorithm}
\usepackage{algorithmic}
\usepackage{listings}
\usepackage{xcolor}

\addbibresource{references.bib}

\geometry{left=2.5cm,right=2.5cm,top=2.5cm,bottom=2.5cm}
\onehalfspacing

% Code listing style
\lstset{
    basicstyle=\ttfamily\footnotesize,
    breaklines=true,
    frame=single,
    language=Python,
    showstringspaces=false,
    commentstyle=\color{gray},
    keywordstyle=\color{blue},
    stringstyle=\color{red}
}

\title{\textbf{Communication Range Effects on Multi-Agent Fruit Harvesting Cooperation: An Agent-Based Modeling Study}}
\author{Ahmed Maruf \\ MSc in Artificial Intelligence \\ Student ID: 250046920}
\date{\today}

\begin{document}

\maketitle

\begin{abstract}
Multi-agent systems (MAS) in resource gathering scenarios face a fundamental trade-off in communication design: insufficient communication leads to poor coordination and missed opportunities, while excessive communication can cause congestion, herding behaviors, and reduced overall performance. This dissertation investigates the effects of communication range on cooperative behavior in multi-agent fruit harvesting through agent-based modeling using the Mesa framework. We develop a spatially explicit simulation where autonomous harvesting agents can share information about resource locations within a limited communication range, implementing both cooperative and competitive strategies. Our research hypothesis posits that there exists an optimal communication range that maximizes collective harvesting efficiency by balancing information sharing benefits against coordination overhead and spatial herding effects. Through systematic experimentation across varying communication ranges (0-8 grid cells), agent densities (10-30 agents), and resource distributions (uniform vs. clustered), we aim to identify this optimal range and understand the underlying mechanisms. Preliminary implementation includes a Mesa-based simulation with HarvesterAgent classes featuring position tracking, fruit inventory management, range-bounded message delivery, and interactive web visualization. The study contributes to understanding communication-cooperation dynamics in distributed autonomous systems with applications to swarm robotics, distributed sensing networks, and collaborative AI systems. Expected outcomes include identification of optimal communication ranges, validation of the cooperation-communication trade-off hypothesis, and insights into scalable coordination mechanisms for multi-agent resource gathering tasks.
\end{abstract}

\section{Introduction}
\label{sec:introduction}

The coordination of multiple autonomous agents in resource gathering tasks represents a fundamental challenge in distributed artificial intelligence and multi-agent systems. From swarm robotics applications in environmental monitoring to distributed sensor networks in precision agriculture, the ability of agents to effectively share information and coordinate their actions directly impacts system performance and efficiency \parencite{dorigo2004ant, cao1997cooperative}.

A critical design parameter in such systems is the communication range—the maximum distance over which agents can directly exchange information. This parameter creates a fundamental trade-off: too little communication leads to poor coordination, redundant exploration, and missed opportunities for collaboration, while excessive communication can result in information overload, network congestion, and herding behaviors that reduce spatial coverage and overall system efficiency \parencite{reynolds1987flocks, vicsek1995novel}.

\subsection{Research Motivation}

The motivation for this research stems from both theoretical gaps in multi-agent systems literature and practical challenges in deploying cooperative autonomous systems. While extensive research exists on communication protocols and coordination mechanisms in multi-agent systems \parencite{stone2000multiagent, tambe1997towards}, there is limited systematic investigation into the optimal communication range for resource gathering tasks, particularly in spatially distributed environments where agents must balance local coordination with global exploration.

Real-world applications highlight the importance of this research question. In precision agriculture, autonomous harvesting robots must coordinate to efficiently cover orchards while avoiding interference \parencite{bac2014performance}. In environmental monitoring, sensor networks must balance local information sharing with network-wide coverage \parencite{akyildiz2002wireless}. In search and rescue operations, robot swarms must coordinate search patterns while maintaining communication links \parencite{murphy2008human}.

\subsection{Research Questions and Objectives}

This dissertation addresses the following primary research question:

\textbf{Primary Research Question:} What is the optimal communication range for multi-agent fruit harvesting cooperation, and how does this range vary with agent density, resource distribution, and cooperation strategy?

\textbf{Supporting Research Questions:}
\begin{enumerate}
    \item \textbf{RQ1: Communication-Performance Relationship} - How does communication range affect collective harvesting performance, and is there evidence for an optimal intermediate range?
    \item \textbf{RQ2: Cooperation vs. Competition} - How do cooperative and competitive agent strategies interact with communication range to influence system performance?
    \item \textbf{RQ3: Scalability and Density Effects} - How does the optimal communication range change with agent density and spatial distribution?
    \item \textbf{RQ4: Resource Distribution Impact} - How do different resource spatial patterns (uniform vs. clustered) affect the optimal communication range?
\end{enumerate}

\subsection{Research Hypothesis}

We hypothesize that there exists an optimal communication range that maximizes collective harvesting efficiency by balancing the benefits of information sharing against the costs of coordination overhead and spatial herding effects. Specifically:

\textbf{H1:} Performance as a function of communication range follows an inverted-U curve, with peak performance at an intermediate range.

\textbf{H2:} The optimal communication range increases with resource clustering and decreases with agent density.

\textbf{H3:} Cooperative strategies benefit more from communication than competitive strategies, but are also more susceptible to herding effects at large communication ranges.

\subsection{Methodology Overview}

This research employs agent-based modeling using the Mesa framework to create a controlled experimental environment for studying communication range effects. The simulation features autonomous HarvesterAgent entities operating in a discrete grid environment with spatially distributed fruit resources. Agents can share information about resource locations within their communication range and employ either cooperative or competitive harvesting strategies.

The experimental design involves systematic parameter sweeps across communication range (0-8 grid cells), agent density (10-30 agents), and resource distribution patterns, with multiple replications to ensure statistical validity. Key performance metrics include collective harvesting efficiency, spatial coverage, communication overhead, and individual agent performance variance.

\subsection{Contributions and Significance}

This research makes several contributions to the field of multi-agent systems:

\begin{enumerate}
    \item \textbf{Empirical Evidence:} Systematic investigation of communication range effects on cooperative behavior in resource gathering tasks.
    \item \textbf{Methodological Framework:} A reusable Mesa-based simulation framework for studying communication-cooperation dynamics.
    \item \textbf{Design Guidelines:} Practical insights for optimal communication range selection in distributed autonomous systems.
    \item \textbf{Theoretical Understanding:} Enhanced understanding of the trade-offs between information sharing benefits and coordination costs.
\end{enumerate}

The findings have practical applications in swarm robotics, distributed sensor networks, autonomous vehicle coordination, and collaborative AI systems where communication range is a critical design parameter.

\subsection{Dissertation Structure}

The remainder of this dissertation is organized as follows: Section \ref{sec:literature} reviews relevant literature on multi-agent coordination, communication in distributed systems, and swarm intelligence. Section \ref{sec:methodology} describes the agent-based modeling approach, simulation design, and experimental methodology. Section \ref{sec:implementation} details the Mesa-based implementation including agent architecture and visualization components. Section \ref{sec:results} presents experimental results and statistical analysis. Section \ref{sec:discussion} discusses findings, implications, and limitations. Section \ref{sec:conclusion} concludes with contributions and future work directions.

\section{Literature Review}
\label{sec:literature}

This section reviews the relevant literature across three key areas: multi-agent coordination and cooperation, communication range optimization in distributed systems, and swarm intelligence approaches to collective behavior. We identify gaps in current research and position our work within the broader context of multi-agent systems research.

\subsection{Multi-Agent Coordination and Cooperation}

Multi-agent coordination has been a central research area since the emergence of distributed artificial intelligence \parencite{bond1988readings}. Early work by \textcite{smith1980contract} introduced the Contract Net Protocol, establishing foundational principles for task allocation in multi-agent systems. \textcite{tambe1997towards} extended this work by developing teamwork models that emphasize shared mental models and coordinated action selection.

In resource gathering scenarios, coordination becomes particularly challenging due to the spatial and temporal distribution of resources. \textcite{stone2000multiagent} demonstrated that effective coordination requires balancing individual agent autonomy with collective objectives. Their work on RoboCup soccer highlighted how communication constraints affect team performance, showing that limited communication can sometimes improve performance by reducing coordination overhead.

Recent work by \textcite{jarne2022mean} provides a mean-field analysis of collaborative multi-agent foragers, showing how coupling strength between agents affects collective performance. Their mathematical framework demonstrates that increased communication can lead to synchronization effects that may harm overall system performance—a key insight supporting our hypothesis of an optimal intermediate communication range.

\textcite{soma2023congestion} specifically studied congestion effects in robot swarms, demonstrating that both movement and communication congestion can drastically affect scalability. Their work showed that "versioned local communication" scales better than naive broadcasting, providing empirical support for the importance of communication range optimization.

\subsection{Communication in Multi-Agent Systems}

Communication in multi-agent systems has evolved from simple message passing to sophisticated protocols that consider bandwidth limitations, network topology, and information relevance \parencite{finin1994kqml, labrou1997agent}. \textcite{goldman2003netlogo} showed that communication range significantly affects the emergence of cooperative behaviors in spatial environments.

The trade-off between communication benefits and costs has been explored in various contexts. \textcite{foerster2016learning} demonstrated that agents can learn when and what to communicate in cooperative multi-agent reinforcement learning settings. Their work showed that selective communication often outperforms full communication, supporting the hypothesis that more communication is not always better.

\textcite{singh2019ic3net} extended this work by developing gated communication mechanisms that allow agents to learn when to communicate. Their results showed that targeted communication based on local observations can achieve better performance than broadcast communication, particularly in resource-constrained environments.

In swarm robotics, \textcite{jimenez2023emergent} studied how communication enhances foraging behavior in evolved swarms controlled by spiking neural networks. They found that emergent communication improves performance but noted that overly diffuse signaling can homogenize behavior, reducing the diversity needed for effective exploration.

\subsection{Swarm Intelligence and Collective Behavior}

Swarm intelligence research has provided insights into how simple local interactions can lead to complex collective behaviors \parencite{bonabeau1999swarm, kennedy2001swarm}. In foraging contexts, \textcite{deneubourg1990self} showed how ant colonies use stigmergic communication (pheromone trails) to coordinate foraging activities without direct agent-to-agent communication.

\textcite{hecker2015beyond} studied ant-inspired robot swarms and found that minimal local rules can yield scalable cooperation. Their work emphasized that locality in communication is often sufficient and more robust than broad signaling, supporting the search for optimal intermediate communication ranges.

The Moses Lab at University of New Mexico has conducted extensive research on swarm foraging. \textcite{lu2020swarm} provided a comprehensive review of swarm foraging algorithms, highlighting the gap between theoretical proof and practical implementation. They emphasized the importance of considering communication overhead and congestion effects in real-world deployments.

\textcite{flanagan2012quantifying} studied the effects of colony size and food distribution on harvester ant foraging, showing that resource clustering interacts with social information sharing. When resources are clustered, local communication helps; when resources are uniform, additional communication provides diminishing returns.

\subsection{Agent-Based Modeling Frameworks}

Agent-based modeling has become a standard approach for studying multi-agent systems \parencite{wilensky2015introduction}. The Mesa framework \parencite{masad2015mesa} provides a Python-based platform for developing agent-based models with built-in visualization and data collection capabilities.

\textcite{pinciroli2012argos} developed ARGoS, a physics-based simulator for multi-robot systems that has become a standard for swarm robotics research. Their work established methodological best practices for large-scale swarm experiments, including the importance of multiple random seeds and statistical validation.

Recent work on explainable agent-based models \parencite{barnes2021explaining} introduced principled simplification as a method for understanding complex multi-agent behaviors. Their approach of creating minimal testbeds that preserve essential problem structure while removing confounding factors provides a methodological framework that we adopt in our experimental design.

\subsection{Research Gaps and Positioning}

Despite extensive research in multi-agent coordination and swarm intelligence, several gaps remain:

\begin{enumerate}
    \item \textbf{Systematic Communication Range Analysis:} While many studies acknowledge the importance of communication range, few provide systematic empirical analysis of optimal ranges across different conditions.

    \item \textbf{Cooperation-Competition Dynamics:} Limited research exists on how communication range affects the balance between cooperative and competitive behaviors in resource gathering tasks.

    \item \textbf{Scalability Considerations:} Most studies focus on small-scale systems; the scalability of communication range effects to larger agent populations remains understudied.

    \item \textbf{Practical Design Guidelines:} There is a lack of practical guidelines for selecting communication ranges in real-world multi-agent systems.
\end{enumerate}

This dissertation addresses these gaps by providing systematic empirical analysis of communication range effects in multi-agent fruit harvesting, with specific attention to cooperation-competition dynamics, scalability considerations, and practical design implications.

Our work builds on the theoretical foundations established by \textcite{jarne2022mean} and \textcite{soma2023congestion}, the methodological approaches of \textcite{barnes2021explaining}, and the practical insights from swarm foraging research \parencite{lu2020swarm, hecker2015beyond}. By combining these perspectives in a systematic experimental framework, we aim to provide both theoretical understanding and practical guidance for communication range optimization in multi-agent systems.

\section{Methodology}
\label{sec:methodology}

This section describes our research methodology, including the agent-based modeling approach, simulation design, experimental framework, and data analysis methods. We employ a quantitative experimental design using controlled simulations to systematically investigate communication range effects on multi-agent cooperation.

\subsection{Research Design Overview}

Our research follows a quantitative experimental approach using agent-based modeling to create a controlled environment for studying communication range effects. The methodology is designed to address our research questions through systematic parameter manipulation and statistical analysis of resulting behaviors.

The experimental design employs a full factorial approach across key parameters: communication range, agent density, resource distribution, and cooperation strategy. This allows us to identify main effects and interactions between these factors while maintaining statistical rigor through multiple replications and appropriate controls.

\subsection{Agent-Based Modeling Framework}

We selected agent-based modeling as our primary methodology for several reasons:

\begin{enumerate}
    \item \textbf{Spatial Representation:} ABM naturally represents spatial relationships crucial for communication range effects.
    \item \textbf{Emergent Behavior:} ABM captures emergent collective behaviors arising from individual agent interactions.
    \item \textbf{Parameter Control:} ABM allows precise control over experimental parameters while maintaining system complexity.
    \item \textbf{Scalability:} ABM enables testing across different scales and agent populations.
\end{enumerate}

\subsubsection{Mesa Framework Selection}

We chose the Mesa framework \parencite{masad2015mesa} for implementation based on the following criteria:

\begin{itemize}
    \item \textbf{Python Integration:} Seamless integration with scientific Python ecosystem (NumPy, Pandas, Matplotlib)
    \item \textbf{Built-in Visualization:} Interactive web-based visualization for model exploration
    \item \textbf{Data Collection:} Integrated DataCollector for systematic metric tracking
    \item \textbf{Grid Spaces:} Native support for discrete grid environments with neighborhood operations
    \item \textbf{Reproducibility:} Built-in random seed management for reproducible experiments
\end{itemize}

\subsection{Simulation Environment Design}

\subsubsection{Spatial Environment}

The simulation environment consists of a discrete 2D grid representing an orchard or foraging area. Key environmental parameters include:

\begin{itemize}
    \item \textbf{Grid Size:} 30×30 cells (scalable for density experiments)
    \item \textbf{Boundary Conditions:} Non-toroidal (bounded) to simulate realistic spatial constraints
    \item \textbf{Cell Occupancy:} Multiple agents can occupy the same cell (realistic for fruit harvesting)
    \item \textbf{Neighborhood:} Moore neighborhood (8-connected) for movement and communication
\end{itemize}

\subsubsection{Resource Distribution}

Fruit resources are distributed according to two primary patterns:

\begin{enumerate}
    \item \textbf{Uniform Distribution:} Random placement with equal probability across all cells
    \item \textbf{Clustered Distribution:} Gaussian clusters with varying cluster sizes and densities
\end{enumerate}

Resource dynamics include:
\begin{itemize}
    \item \textbf{Initial Density:} 0.3-0.5 fruits per cell on average
    \item \textbf{Regeneration:} Probabilistic regeneration (0.1 probability per time step per empty cell)
    \item \textbf{Depletion:} Fruits are consumed upon harvesting
\end{itemize}

\subsection{Agent Architecture}

\subsubsection{HarvesterAgent Design}

Each HarvesterAgent is implemented as a Mesa Agent with the following core components:

\begin{algorithm}
\caption{HarvesterAgent Step Procedure}
\begin{algorithmic}[1]
\STATE \textbf{Sense:} Detect local fruit and receive messages
\STATE \textbf{Decide:} Select target location and action
\STATE \textbf{Communicate:} Send messages within communication range
\STATE \textbf{Act:} Move toward target and harvest if possible
\STATE \textbf{Update:} Update internal state and memory
\end{algorithmic}
\end{algorithm}

\textbf{Agent State Variables:}
\begin{itemize}
    \item \textbf{Position:} Current grid coordinates (x, y)
    \item \textbf{Communication Range:} Maximum distance for message transmission (R)
    \item \textbf{Strategy:} Cooperation mode ("cooperative" or "competitive")
    \item \textbf{Fruit Inventory:} Count of harvested fruits
    \item \textbf{Known Locations:} Memory of fruit locations from direct observation and communication
    \item \textbf{Message Inbox:} Buffer for received messages
\end{itemize}

\subsubsection{Communication Model}

The communication system implements range-bounded message delivery with the following characteristics:

\begin{itemize}
    \item \textbf{Range Constraint:} Messages only reach agents within Euclidean distance R
    \item \textbf{Message Types:} HOTSPOT (fruit location), CLAIM (target reservation), RELEASE (target abandonment)
    \item \textbf{Broadcast Model:} All agents within range receive all messages (no selective addressing)
    \item \textbf{Instantaneous Delivery:} No communication delays or failures (idealized model)
\end{itemize}

\subsubsection{Cooperation Strategies}

Two primary strategies are implemented:

\begin{enumerate}
    \item \textbf{Cooperative Strategy:}
    \begin{itemize}
        \item Shares fruit location information via HOTSPOT messages
        \item Respects CLAIM messages from other agents
        \item Avoids targeting claimed resources
    \end{itemize}

    \item \textbf{Competitive Strategy:}
    \begin{itemize}
        \item Does not share fruit location information
        \item Ignores CLAIM messages from other agents
        \item Pursues any detected fruit regardless of claims
    \end{itemize}
\end{enumerate}

\subsection{Experimental Design}

\subsubsection{Independent Variables}

Our experimental design manipulates four primary independent variables:

\begin{enumerate}
    \item \textbf{Communication Range (R):} 0, 1, 2, 3, 4, 6, 8 grid cells
    \item \textbf{Agent Density:} 10, 20, 30 agents (on 30×30 grid)
    \item \textbf{Resource Distribution:} Uniform vs. Clustered
    \item \textbf{Cooperation Strategy:} 100\% Cooperative, 100\% Competitive, 50\% Mixed
\end{enumerate}

\subsubsection{Dependent Variables}

Key performance metrics include:

\begin{itemize}
    \item \textbf{Collective Efficiency:} Total fruits harvested per time step
    \item \textbf{Individual Performance:} Fruits harvested per agent per time step
    \item \textbf{Spatial Coverage:} Unique cells visited by all agents
    \item \textbf{Communication Overhead:} Messages sent per agent per time step
    \item \textbf{Coordination Conflicts:} Instances of multiple agents targeting the same fruit
    \item \textbf{Network Connectivity:} Size of largest connected component in communication graph
\end{itemize}

\subsubsection{Experimental Conditions}

The full experimental design includes:
\begin{itemize}
    \item \textbf{Total Conditions:} 7 (ranges) × 3 (densities) × 2 (distributions) × 3 (strategies) = 126 conditions
    \item \textbf{Replications:} 30 independent runs per condition
    \item \textbf{Episode Length:} 200 time steps per run
    \item \textbf{Total Simulations:} 126 × 30 = 3,780 simulation runs
\end{itemize}

\subsection{Data Collection and Analysis}

\subsubsection{Data Collection Framework}

Mesa's DataCollector is configured to record both agent-level and model-level metrics at each time step:

\begin{lstlisting}[language=Python, caption=Data Collection Configuration]
datacollector = DataCollector(
    model_reporters={
        "Total_Fruits_Harvested": lambda m: sum(a.fruits_collected
                                               for a in m.agents),
        "Messages_Sent": lambda m: m.total_messages_sent,
        "Spatial_Coverage": compute_spatial_coverage,
        "Network_Connectivity": compute_network_connectivity
    },
    agent_reporters={
        "Fruits_Collected": "fruits_collected",
        "Messages_Sent": "messages_sent",
        "Position": "pos"
    }
)
\end{lstlisting}

\subsubsection{Statistical Analysis Plan}

Our analysis employs multiple statistical approaches:

\begin{enumerate}
    \item \textbf{Descriptive Statistics:} Mean, standard deviation, and confidence intervals for all metrics
    \item \textbf{ANOVA:} Multi-factor analysis of variance to identify main effects and interactions
    \item \textbf{Regression Analysis:} Polynomial regression to identify optimal communication ranges
    \item \textbf{Post-hoc Testing:} Tukey's HSD for pairwise comparisons between communication ranges
    \item \textbf{Effect Size:} Cohen's d and eta-squared to quantify practical significance
\end{enumerate}

\subsubsection{Hypothesis Testing}

Specific statistical tests for each hypothesis:

\begin{itemize}
    \item \textbf{H1 (Inverted-U Curve):} Quadratic regression analysis with significance testing for quadratic term
    \item \textbf{H2 (Density and Clustering Effects):} Two-way ANOVA with interaction terms
    \item \textbf{H3 (Strategy Differences):} Mixed-effects ANOVA with strategy as between-subjects factor
\end{itemize}

\subsection{Validation and Verification}

\subsubsection{Model Verification}

We employ several verification techniques:

\begin{enumerate}
    \item \textbf{Unit Testing:} Individual agent behaviors and communication functions
    \item \textbf{Boundary Testing:} Extreme parameter values (R=0, R=max)
    \item \textbf{Conservation Laws:} Total fruit conservation (created = harvested + remaining)
    \item \textbf{Deterministic Testing:} Fixed random seeds for reproducible behavior
\end{enumerate}

\subsubsection{Model Validation}

Validation approaches include:

\begin{enumerate}
    \item \textbf{Face Validity:} Expert review of agent behaviors and emergent patterns
    \item \textbf{Sensitivity Analysis:} Parameter perturbation to assess model robustness
    \item \textbf{Extreme Conditions:} Testing behavior at parameter boundaries
    \item \textbf{Literature Comparison:} Qualitative comparison with published swarm foraging results
\end{enumerate}

\subsection{Ethical Considerations}

This research involves computational simulations with no human subjects or sensitive data. However, we consider broader ethical implications:

\begin{itemize}
    \item \textbf{Reproducibility:} All code and data will be made publicly available
    \item \textbf{Transparency:} Clear documentation of assumptions and limitations
    \item \textbf{Responsible AI:} Consideration of potential applications in autonomous systems
\end{itemize}

The methodology described above provides a rigorous framework for investigating communication range effects in multi-agent cooperation, with appropriate controls, statistical power, and validation procedures to ensure reliable and meaningful results.

\section{Implementation}
\label{sec:implementation}

This section details the technical implementation of our multi-agent fruit harvesting simulation using the Mesa framework. We describe the software architecture, agent implementation, communication system, visualization components, and current development status according to our 12-week implementation plan.

\subsection{Software Architecture}

\subsubsection{Project Structure}

The simulation is organized into a modular Python package structure:

\begin{lstlisting}[language=bash, caption=Project Directory Structure]
fruit_harvester_simulation/
├── __init__.py              # Package initialization
├── agents.py                # Agent class definitions
├── environment.py           # Environment and model classes
├── simulation.py            # Simulation runner and batch processing
├── viz_server.py           # Mesa web visualization server
└── data_collection.py      # Data analysis utilities (planned)
\end{lstlisting}

\subsubsection{Core Dependencies}

The implementation relies on the following key libraries:

\begin{itemize}
    \item \textbf{Mesa 2.4.0:} Agent-based modeling framework
    \item \textbf{Pandas 2.3.0:} Data manipulation and analysis
    \item \textbf{Matplotlib 3.9.4:} Visualization and plotting
    \item \textbf{NumPy:} Numerical computations
    \item \textbf{NetworkX:} Graph analysis for communication networks (planned)
\end{itemize}

\subsection{Agent Implementation}

\subsubsection{HarvesterAgent Class}

The core HarvesterAgent class implements the agent architecture described in our methodology:

\begin{lstlisting}[language=Python, caption=HarvesterAgent Core Structure]
class HarvesterAgent(Agent):
    def __init__(self, unique_id, model, communication_range,
                 strategy="cooperative", sensor_radius=2):
        super().__init__(unique_id, model)
        self.communication_range = communication_range
        self.strategy = strategy
        self.sensor_radius = sensor_radius
        self.fruits_collected = 0
        self.known_fruit_locations = set()
        self.inbox = []
        self.outbox = []

    def step(self):
        """Main agent step following sense-decide-act cycle"""
        self.sense()
        self.decide()
        self.act()
        self.communicate()
\end{lstlisting}

\subsubsection{Sensing and Perception}

Agents perceive their environment through local sensing and message reception:

\begin{lstlisting}[language=Python, caption=Agent Sensing Implementation]
def sense(self):
    """Detect nearby fruits and process incoming messages"""
    # Local fruit detection
    nearby_cells = self.model.grid.get_neighborhood(
        self.pos, moore=True, radius=self.sensor_radius
    )

    for cell in nearby_cells:
        fruit_count = self.model.get_fruit_at(cell)
        if fruit_count > 0:
            self.known_fruit_locations.add(cell)

    # Process incoming messages
    self._process_inbox()
\end{lstlisting}

\subsubsection{Decision Making}

The decision-making process selects target locations based on known fruit locations and strategy:

\begin{lstlisting}[language=Python, caption=Agent Decision Making]
def decide(self):
    """Select target location and action"""
    if not self.known_fruit_locations:
        # Random exploration if no known fruits
        self.target = self._select_random_target()
    else:
        # Select closest known fruit location
        self.target = min(self.known_fruit_locations,
                         key=lambda loc: self._distance_to(loc))

        # Cooperative agents check for claims
        if self.strategy == "cooperative":
            if self._is_claimed(self.target):
                self.target = self._select_alternative_target()
\end{lstlisting}

\subsection{Communication System}

\subsubsection{Message Types}

The communication system implements three message types:

\begin{lstlisting}[language=Python, caption=Message Class Definition]
class Message:
    def __init__(self, sender_id, message_type, content, ttl=5):
        self.sender_id = sender_id
        self.message_type = message_type  # HOTSPOT, CLAIM, RELEASE
        self.content = content
        self.ttl = ttl  # Time-to-live for message expiration
        self.timestamp = 0
\end{lstlisting}

\subsubsection{Range-Bounded Message Delivery}

The OrchardModel implements range-bounded message delivery:

\begin{lstlisting}[language=Python, caption=Message Delivery System]
def deliver_messages(self):
    """Deliver messages within communication range"""
    for sender in self.agents:
        if not sender.outbox:
            continue

        # Find agents within communication range
        nearby_agents = []
        for receiver in self.agents:
            if sender.unique_id != receiver.unique_id:
                distance = self._euclidean_distance(
                    sender.pos, receiver.pos
                )
                if distance <= sender.communication_range:
                    nearby_agents.append(receiver)

        # Deliver messages to nearby agents
        for message in sender.outbox:
            for receiver in nearby_agents:
                receiver.inbox.append(message)

        sender.outbox.clear()
\end{lstlisting}

\subsection{Environment Implementation}

\subsubsection{OrchardModel Class}

The main simulation model extends Mesa's Model class:

\begin{lstlisting}[language=Python, caption=OrchardModel Structure]
class OrchardModel(Model):
    def __init__(self, width=30, height=30, num_agents=20,
                 communication_range=3, cooperative_share=1.0):
        super().__init__()
        self.grid = MultiGrid(width, height, torus=False)
        self.schedule = RandomActivation(self)

        # Initialize fruit map
        self.fruit_map = np.zeros((width, height))
        self._initial_fruit_spawn()

        # Create agents
        self._create_agents(num_agents, communication_range,
                          cooperative_share)

        # Data collection
        self.datacollector = DataCollector(
            model_reporters=self._get_model_reporters(),
            agent_reporters=self._get_agent_reporters()
        )
\end{lstlisting}

\subsubsection{Fruit Dynamics}

The environment manages fruit spawning, regeneration, and harvesting:

\begin{lstlisting}[language=Python, caption=Fruit Management System]
def regenerate_fruit(self):
    """Probabilistic fruit regeneration"""
    for x in range(self.grid.width):
        for y in range(self.grid.height):
            if self.fruit_map[x, y] == 0:  # Empty cell
                if self.random.random() < self.regeneration_prob:
                    self.fruit_map[x, y] = 1

def harvest_fruit(self, pos):
    """Harvest fruit at given position"""
    x, y = pos
    if self.fruit_map[x, y] > 0:
        self.fruit_map[x, y] -= 1
        return True
    return False
\end{lstlisting}

\subsection{Visualization System}

\subsubsection{Web-Based Visualization}

We implemented an interactive web visualization using Mesa's visualization framework:

\begin{lstlisting}[language=Python, caption=Visualization Server Setup]
def launch_visualization():
    """Launch Mesa web visualization server"""
    grid = CanvasGrid(portrayal_method, 30, 30, 750, 750)

    chart = ChartModule([
        {"Label": "Total Fruits Collected", "Color": "Blue"},
        {"Label": "Messages Sent", "Color": "Red"}
    ])

    server = ModularServer(
        OrchardModel,
        [grid, chart],
        "Fruit Harvesting Simulation",
        model_params
    )

    server.port = 8521
    server.launch()
\end{lstlisting}

\subsubsection{Agent and Environment Portrayal}

The visualization system renders agents and environment elements:

\begin{lstlisting}[language=Python, caption=Visualization Portrayal]
def portrayal_method(agent):
    """Define visual representation of agents and environment"""
    if isinstance(agent, HarvesterAgent):
        color = "#00CED1" if agent.strategy == "cooperative" else "#FF00FF"
        return {
            "Shape": "circle",
            "r": 0.5,
            "Filled": "true",
            "Color": color,
            "Layer": 1,
            "text": str(agent.unique_id)
        }
    elif isinstance(agent, CellPatch):
        # Background showing fruit density
        fruit_count = agent.model.get_fruit_at(agent.pos)
        intensity = min(fruit_count / 3.0, 1.0)
        color_value = int(255 * (1 - intensity * 0.3))
        color = f"rgb(255, {color_value}, {color_value})"
        return {
            "Shape": "rect",
            "w": 1, "h": 1,
            "Color": color,
            "Layer": 0
        }
\end{lstlisting}

\subsection{Current Implementation Status}

\subsubsection{Completed Components (Week 2 Status)}

According to our 12-week implementation plan, we have completed the Week 2 deliverables:

\begin{itemize}
    \item ✅ \textbf{HarvesterAgent Class:} Core agent with position tracking, fruit inventory, and communication range
    \item ✅ \textbf{Communication System:} Range-bounded message delivery with HOTSPOT, CLAIM, and RELEASE messages
    \item ✅ \textbf{Cooperation Strategies:} Cooperative vs. competitive agent behaviors
    \item ✅ \textbf{Environment:} Fruit spawning, regeneration, and harvesting mechanics
    \item ✅ \textbf{Visualization:} Interactive web-based visualization with agent and fruit display
    \item ✅ \textbf{Basic Testing:} Smoke tests confirming simulation runs successfully
\end{itemize}

\subsubsection{Week 3 Planned Components}

The next implementation phase (Week 3) focuses on data collection and metrics:

\begin{itemize}
    \item 🔄 \textbf{DataCollector Integration:} Comprehensive metric tracking
    \item 🔄 \textbf{Performance Metrics:} Collective efficiency, spatial coverage, communication overhead
    \item 🔄 \textbf{Data Export:} CSV export functionality for statistical analysis
    \item 🔄 \textbf{Batch Processing:} Automated experiment execution across parameter ranges
\end{itemize}

\subsubsection{Technical Challenges and Solutions}

Several technical challenges were encountered and resolved:

\begin{enumerate}
    \item \textbf{Grid Coordinate Systems:} Mesa's coordinate system required careful handling of (x,y) vs. (row,col) conventions.

    \item \textbf{Message Delivery Timing:} Implemented staged execution (sense → decide → communicate → act) to avoid race conditions.

    \item \textbf{Visualization Performance:} Used CellPatch agents for background rendering to improve visualization responsiveness.

    \item \textbf{Random Number Generation:} Ensured reproducible results through proper random seed management.
\end{enumerate}

\subsection{Code Quality and Documentation}

\subsubsection{Testing Framework}

Basic testing infrastructure is in place:

\begin{lstlisting}[language=Python, caption=Basic Test Structure]
def test_agent_creation():
    """Test agent initialization and basic properties"""
    model = OrchardModel(width=10, height=10, num_agents=5)
    assert len(model.agents) == 5
    assert all(isinstance(agent, HarvesterAgent) for agent in model.agents)

def test_communication_range():
    """Test message delivery within communication range"""
    model = OrchardModel(width=10, height=10, num_agents=2,
                        communication_range=3)
    # Place agents at known positions and test message delivery
    # ... test implementation
\end{lstlisting}

\subsubsection{Documentation Standards}

Code follows Python documentation standards with docstrings for all classes and methods:

\begin{lstlisting}[language=Python, caption=Documentation Example]
def _euclidean_distance(self, pos1, pos2):
    """
    Calculate Euclidean distance between two grid positions.

    Args:
        pos1 (tuple): First position (x, y)
        pos2 (tuple): Second position (x, y)

    Returns:
        float: Euclidean distance between positions
    """
    return math.sqrt((pos1[0] - pos2[0])**2 + (pos1[1] - pos2[1])**2)
\end{lstlisting}

\subsection{Performance Considerations}

\subsubsection{Computational Complexity}

Current implementation complexity analysis:

\begin{itemize}
    \item \textbf{Message Delivery:} O(n²) per step for n agents (all-pairs distance calculation)
    \item \textbf{Fruit Detection:} O(k) per agent for k cells in sensor radius
    \item \textbf{Pathfinding:} O(1) simple movement toward target
    \item \textbf{Overall:} O(n²) per simulation step
\end{itemize}

\subsubsection{Optimization Opportunities}

Future optimizations for larger-scale experiments:

\begin{enumerate}
    \item \textbf{Spatial Indexing:} Use spatial data structures for efficient neighbor queries
    \item \textbf{Message Filtering:} Implement message relevance filtering to reduce communication overhead
    \item \textbf{Parallel Processing:} Utilize multiprocessing for batch experiment execution
    \item \textbf{Caching:} Cache distance calculations and neighborhood queries
\end{enumerate}

The current implementation provides a solid foundation for our experimental investigation, with all core components functional and ready for the data collection and analysis phases of our research.

\section{Preliminary Results and Planned Experiments}
\label{sec:results}

This section presents preliminary results from initial simulation runs and outlines the comprehensive experimental plan for investigating communication range effects on multi-agent cooperation. While full experimental results are pending completion of the data collection framework (Week 3 of implementation), we present initial observations and the detailed experimental design.

\subsection{Initial Simulation Validation}

\subsubsection{Basic Functionality Testing}

Initial testing confirms that the core simulation components function correctly:

\begin{itemize}
    \item \textbf{Agent Movement:} Agents successfully navigate the grid environment and move toward fruit targets
    \item \textbf{Fruit Harvesting:} Fruit collection and inventory tracking work as designed
    \item \textbf{Communication:} Message delivery operates within specified communication ranges
    \item \textbf{Strategy Differentiation:} Cooperative and competitive agents exhibit distinct behaviors
\end{itemize}

\subsubsection{Smoke Test Results}

A preliminary smoke test with 10 agents, communication range 3, and 10 simulation steps yielded:

\begin{itemize}
    \item \textbf{Total Fruits Collected:} 40 fruits over 10 steps
    \item \textbf{Message Volume:} Cooperative agents generated 15-20 messages per step
    \item \textbf{Spatial Distribution:} Agents spread across the environment without excessive clustering
    \item \textbf{Performance Stability:} Simulation runs consistently without errors
\end{itemize}

\subsection{Planned Experimental Framework}

\subsubsection{Experimental Design Matrix}

The full experimental design will systematically investigate the following parameter combinations:

\begin{table}[h]
\centering
\caption{Experimental Design Parameters}
\label{tab:experimental_design}
\begin{tabular}{@{}lll@{}}
\toprule
\textbf{Parameter} & \textbf{Values} & \textbf{Rationale} \\
\midrule
Communication Range & 0, 1, 2, 3, 4, 6, 8 cells & Cover no-comm to high-comm spectrum \\
Agent Density & 10, 20, 30 agents & Test scalability effects \\
Resource Distribution & Uniform, Clustered & Examine spatial pattern effects \\
Strategy Mix & 100\% Coop, 100\% Comp, 50\% Mixed & Compare cooperation levels \\
Replications & 30 per condition & Ensure statistical power \\
Episode Length & 200 time steps & Allow convergence to steady state \\
\bottomrule
\end{tabular}
\end{table}

\subsubsection{Hypothesis Testing Framework}

Each research hypothesis will be tested using specific experimental conditions and statistical methods:

\textbf{H1: Inverted-U Performance Curve}
\begin{itemize}
    \item \textbf{Test:} Quadratic regression of performance vs. communication range
    \item \textbf{Conditions:} All agent densities and resource distributions
    \item \textbf{Metric:} Collective harvesting efficiency (fruits per time step)
    \item \textbf{Expected Result:} Significant quadratic term with negative coefficient
\end{itemize}

\textbf{H2: Density and Clustering Effects}
\begin{itemize}
    \item \textbf{Test:} Two-way ANOVA with interaction terms
    \item \textbf{Factors:} Communication range × Agent density × Resource distribution
    \item \textbf{Expected Result:} Optimal range increases with clustering, decreases with density
\end{itemize}

\textbf{H3: Strategy-Communication Interactions}
\begin{itemize}
    \item \textbf{Test:} Mixed-effects ANOVA comparing strategy types
    \item \textbf{Conditions:} All communication ranges with pure strategy conditions
    \item \textbf{Expected Result:} Cooperative agents show stronger communication range effects
\end{itemize}

\subsection{Planned Performance Metrics}

\subsubsection{Primary Metrics}

\begin{enumerate}
    \item \textbf{Collective Efficiency:} $E_c = \frac{\text{Total fruits harvested}}{\text{Time steps} \times \text{Number of agents}}$

    \item \textbf{Spatial Coverage:} $S_c = \frac{\text{Unique cells visited}}{\text{Total grid cells}}$

    \item \textbf{Communication Overhead:} $C_o = \frac{\text{Messages sent per step}}{\text{Number of agents}}$

    \item \textbf{Coordination Efficiency:} $C_e = \frac{\text{Successful harvests}}{\text{Harvest attempts}}$
\end{enumerate}

\subsubsection{Secondary Metrics}

\begin{enumerate}
    \item \textbf{Network Connectivity:} Size of largest connected component in communication graph
    \item \textbf{Information Propagation:} Average time for fruit location information to spread
    \item \textbf{Agent Specialization:} Variance in individual agent performance
    \item \textbf{Resource Utilization:} Fraction of available fruits harvested per episode
\end{enumerate}

\subsection{Expected Results and Predictions}

\subsubsection{Communication Range Effects}

Based on literature review and theoretical considerations, we predict:

\begin{itemize}
    \item \textbf{R = 0 (No Communication):} Poor coordination, redundant exploration, moderate performance
    \item \textbf{R = 1-2 (Local Communication):} Improved local coordination, reduced conflicts
    \item \textbf{R = 3-4 (Optimal Range):} Peak performance balancing information sharing and autonomy
    \item \textbf{R = 6-8 (High Communication):} Herding effects, congestion, reduced performance
\end{itemize}

\subsubsection{Interaction Effects}

Predicted interaction patterns:

\begin{enumerate}
    \item \textbf{Density × Communication Range:} Higher density systems will have lower optimal communication ranges due to increased congestion

    \item \textbf{Resource Distribution × Communication Range:} Clustered resources will benefit from higher communication ranges for information sharing

    \item \textbf{Strategy × Communication Range:} Cooperative agents will show more pronounced communication range effects than competitive agents
\end{enumerate}

\subsection{Statistical Analysis Plan}

\subsubsection{Descriptive Analysis}

For each experimental condition, we will compute:

\begin{itemize}
    \item Mean and standard deviation for all metrics
    \item 95\% confidence intervals
    \item Distribution plots and normality tests
    \item Correlation matrices between metrics
\end{itemize}

\subsubsection{Inferential Statistics}

Primary statistical tests:

\begin{enumerate}
    \item \textbf{ANOVA:} Multi-factor analysis with post-hoc comparisons
    \item \textbf{Regression:} Polynomial regression to identify optimal ranges
    \item \textbf{Effect Size:} Cohen's d and eta-squared for practical significance
    \item \textbf{Power Analysis:} Retrospective power analysis to validate sample sizes
\end{enumerate}

\subsubsection{Visualization Plan}

Planned visualizations include:

\begin{itemize}
    \item Performance curves showing efficiency vs. communication range
    \item Heatmaps of agent movement patterns and spatial coverage
    \item Network diagrams of communication connectivity
    \item Box plots comparing strategy performance across conditions
    \item Interaction plots for multi-factor effects
\end{itemize}

\subsection{Implementation Timeline}

\subsubsection{Week 3-4: Data Collection Framework}

\begin{itemize}
    \item Complete Mesa DataCollector integration
    \item Implement all performance metrics
    \item Create batch processing system for automated experiments
    \item Validate data export and analysis pipeline
\end{itemize}

\subsubsection{Week 4-5: Experimental Execution}

\begin{itemize}
    \item Execute full experimental design (3,780 simulation runs)
    \item Monitor for computational issues and optimize as needed
    \item Implement quality control checks for data integrity
    \item Create preliminary visualizations for ongoing analysis
\end{itemize}

\subsubsection{Week 5-6: Analysis and Results}

\begin{itemize}
    \item Complete statistical analysis of all experimental conditions
    \item Generate publication-quality figures and tables
    \item Validate hypotheses and identify unexpected findings
    \item Prepare comprehensive results section
\end{itemize}

\subsection{Anticipated Challenges}

\subsubsection{Computational Challenges}

\begin{itemize}
    \item \textbf{Scale:} 3,780 simulation runs may require computational optimization
    \item \textbf{Memory:} Large datasets from comprehensive data collection
    \item \textbf{Reproducibility:} Ensuring consistent random seed management across runs
\end{itemize}

\subsubsection{Analytical Challenges}

\begin{itemize}
    \item \textbf{Multiple Comparisons:} Controlling for Type I error across many statistical tests
    \item \textbf{Non-linear Effects:} Identifying complex interaction patterns
    \item \textbf{Outlier Detection:} Managing extreme values in performance metrics
\end{itemize}

\subsubsection{Mitigation Strategies}

\begin{itemize}
    \item \textbf{Parallel Processing:} Utilize multiprocessing for batch experiments
    \item \textbf{Progressive Analysis:} Analyze subsets of data to identify early patterns
    \item \textbf{Robust Statistics:} Use non-parametric tests where appropriate
    \item \textbf{Validation:} Cross-validate findings across different parameter subsets
\end{itemize}

The experimental framework described above provides a comprehensive approach to investigating communication range effects in multi-agent cooperation, with appropriate statistical power and controls to ensure reliable and meaningful results. The next phase of implementation will focus on executing these experiments and analyzing the resulting data to test our hypotheses and contribute new insights to the field of multi-agent systems.

\printbibliography

\end{document}